% 中文正文片段；主文件加载 amsmath, amssymb。
% 批次 20260922T050733-f2dbf5c7；报告 20260922T054144-5aee3c91。
\subsection{三角涡量初值下的粗步长行为与时间收敛}

本研究采用 $\Omega=(0,2\pi)^2$ 上的齐次无滑移问题，
空间离散为 $256\times256$ 均匀交错 MAC 网格，
运动黏性系数为 $\nu=0.001$，保存配置记录的定常外力为
$f(x,y)=(\cos x,0)$。初值身份记为
\texttt{weighted\_trig\_vorticity}，与零初值方腔和周期模型均不同。
本次整理使用该历史初值生成的保存记录，不以当前初值工厂重构或替换历史初态；
初值的具体解析配方未由本次误差整理独立恢复，复现初始化还需对应的历史来源。
比较普通 IMEX--SDIRK2 和阻尼参数 $\gamma=5$ 的增量 mr-ccSAV 格式，
终止时间为 $T=4$，观测时刻为 $t=1,2,4$。

两格式分别使用 $\tau_k=2^{-k}$，$k=0,\ldots,8$。
共同参考为普通 SDIRK2 的 $k=11$ 结果，另以 $k=12$ 进行敏感性检查。
以速度差 $e_h$ 的离散 $L^2$ 范数及黏性二次型
$\sqrt{(e_h,K_he_h)_h}$ 衡量误差，后者为离散 $H^1$ 半范数。
相邻层误差均有限且为正时记录观察阶；失败行保留状态和缺失误差，
不把缺失值替换为零，也不跨越失败层连接误差曲线。

保存结果包含明显的粗步长与细步长区分。普通格式在 $k=1,\ldots,5$
的五个运行失败，对应三个观测时刻共十五行缺失误差；
$k=0$ 虽完成到终点，但 $T=4$ 的 $L^2$ 误差约为
$1.68\times10^{21}$，不能仅凭运行状态视作准确解。
SAV 格式在这些粗步长上给出有限误差，但有限或有界不等于准确，
误差也不必随每一次步长细化单调下降。

在 $k=6,7,8$ 的细步长区间，两格式的曲线接近。
例如 $T=4$ 的普通格式 $L^2$ 误差依次约为
$3.07\times10^{-4}$、$7.65\times10^{-5}$ 和 $1.89\times10^{-5}$，
后两层观察阶约为 $2.004$ 和 $2.017$，
$H^1$ 半范数观察阶也接近二阶。跨越粗细区时出现的异常大观察阶
不是格式的形式阶；因此同时保留全区间图与明确标注的细步长放大图，
不以局部曲线替代全局行为或据此给出无条件准确性结论。

所选报告的两个范数及三个时刻共六项参考敏感性检查通过 $5\%$ 阈值，
但这不是严格的参考误差界。更早的未通过参考检查的报告不与本报告混用。
本次仅核对十八个被测运行及两个参考记录的身份和完整性，
从已保存误差表重绘，不独立重算场误差、补算失败轨迹或进行长期稳定性验证。
